% ======================================================================
\section{Scope, conventions and overview}
\label{sec:scope}
% ======================================================================

\subsection{What is described}

MOCASSIN-PDR couples three pieces of physics:
\begin{enumerate}
\item \textbf{MOCASSIN} \citep{Ercolano2003,Ercolano2005,Ercolano2008}: a
  three-dimensional Monte Carlo (MC) code for photoionisation and dust
  radiative transfer on a Cartesian grid of homogeneous cells. It owns the
  radiation field at every frequency, the dust (optical properties,
  temperatures, charge), the ionisation state of every element, the electron
  density and the gas temperature.
\item \textbf{3D-PDR} \citep{Bisbas2012}, vendored under
  \src{source/\allowbreak{}pdr/\allowbreak{}} (upstream commit \texttt{fabcdc6}, imported 2026-08-28;
  \src{source/\allowbreak{}pdr/\allowbreak{}UPSTREAM.md}): the REDUCED 33-species chemical network solved
  to steady state, its photoreaction prescriptions, its \Htwo{} formation rate
  and its escape-probability level-population solver for the coolants. 3D-PDR's
  own treatment of the radiation field (an external field attenuated along
  rays to the cloud boundary) and its own thermal balance are not used.
  MOCASSIN drives the vendored routines through the bridge module
  \src{source/\allowbreak{}pdr\_mod.f90}.
\item \textbf{The \Htwo{} level model} of the \texttt{h2-model} branch
  (\src{source/\allowbreak{}h2\_mod.f90}, driven by \rtn{h2LevelModel} in
  \src{pdr\_mod.f90}): the 302 rovibrational levels of $\mathrm{H_2}\,X$ in
  statistical equilibrium under collisions, quadrupole decay, ultraviolet
  pumping through the excited electronic states, Solomon dissociation and
  formation on grains. It is a re-implementation from the data and the
  conventions of Cloudy's large \Htwo{} model \citep{Shaw2005,Ferland2017,Chatzikos2023}.
  Its Solomon rate is the \Htwo{} photodissociation rate the chemistry uses.
\end{enumerate}
By \emph{molecular content} we mean the abundances of the 33 species in every
cell and, for \Htwo, the populations of its 302 levels. The document follows
the path by which the MC radiation field determines them, and the feedback of
the molecules on the field, the ionisation and the temperatures.

The text follows the source tree at \texttt{\textasciitilde/software/mocassin\_h2}
(git worktree, branch \texttt{h2-model}, HEAD
\texttt{8fc15120dea2}, 2026-10-15). Every \src{file:\allowbreak{}line} reference is to that
commit; files are under \src{source/\allowbreak{}} unless a path says otherwise. The
branch has moved on since (its tip was \texttt{78a74bd} on 2026-10-16). The only
source change on the path described here is \dlog{D-414} (commit \texttt{df9a626},
merged into \texttt{h2-model} as \texttt{901f1e8}), which makes the mean grain
temperature a mean over the size bins that hold grains
(Section~\ref{sec:Td}); it changes nothing where every grain species fills the
size grid, as in all the fixtures. \dlog{D-417} registers a temperature per
grain material for \Htwo{} formation for future development, without a code
change. DEVLOG
entries are cited as \dlog{D-nnn} (\src{DEVLOG.md}), branch entries as
\dlog{H-nnn} (\src{DEVLOG-H2.md}) and the limitations of
\src{doc/\allowbreak{}pdr-status.md} as \dlog{L-nn}. Measured numbers quoted without a
log entry are read from the run logs of the two production fixtures:
\texttt{hiipdr}, an H\,\textsc{ii} region and PDR around a 40\,kK star in gas of
$n_{\mathrm H}=100$\,cm$^{-3}$ with silicate grains on $31^3$ cells
(\src{runs/\allowbreak{}d405-hii31-d0/\allowbreak{}run.log}), and \texttt{pn150k}, a 150\,kK planetary
nebula with $n_{\mathrm H}=10^3$\,cm$^{-3}$ and graphite grains on $41^3$ cells
(\src{runs/\allowbreak{}d405-pn41-d0/\allowbreak{}run.log}). Both are uniform spheres around one
source, computed in one octant (\kw{symmetricXYZ}).

\paragraph{Build flavour.} The description is of the ray-traced build,
\begin{center}\verb|make PDR=1 PDRGEOM=3D mocassin|\end{center}
which compiles with the flags
\begin{center}\small\ttfamily
-DPDR -DRAYTHEIA -DRAYTHEIA\_MO\\
-DTHERMALBALANCE -DREDUCED -DH2FORM\_CT02 -DCHEMSTEADY -DCHEMRETRY\\
-DFORCECONVERGENCE -DNGACCEL -DNGCYCLE -DNGRELAX -DILLINOIS -DREBRACKET -DSOBOLEV
\end{center}
(\src{Makefile:\allowbreak{}82--89}): the first line selects the ray-traced
geometry, the others are fixed for every PDR build. Of the fixed flags,
\kw{-DREDUCED}, \kw{-DH2FORM\_CT02}, \kw{-DCHEMSTEADY}, \kw{-DCHEMRETRY} and
\kw{-DSOBOLEV} change what the coupled code computes (the network, the \Htwo{}
formation rate, the steady-state integration and the escape-probability
form); \kw{-DTHERMALBALANCE}, \kw{-DFORCECONVERGENCE}, \kw{-DNGACCEL},
\kw{-DNGCYCLE}, \kw{-DNGRELAX}, \kw{-DILLINOIS} and \kw{-DREBRACKET} act only in
3D-PDR's own iteration and temperature routines, which the bridge does not
call. Two other PDR flavours exist and are not
described except where noted: the 1D-slab build (\kw{-DONEDIMENSIONAL}),
whose chemistry still receives 3D-PDR's prescription
$\chi_0\,e^{-\gamma \AV}$ on twelve rays and in which the level model is a
stub that leaves the network's \Htwo{} rate in place
(\src{pdr\_mod.f90:\allowbreak{}6015--6027}); and the warm start
(\verb|make PDR=1 mocassinWarm|). The level model also skips plane-parallel
illumination and nested grids (\src{pdr\_mod.f90:\allowbreak{}5256--5259}), so the
description applies to one grid illuminated by point sources.

\subsection{Notation and units}

Table~\ref{tab:notation} lists the symbols used in more than one section.
cgs units are used throughout; frequencies are in rydbergs in the code
(1\,Ryd $=13.6057$\,eV, $\nu_{1\,\mathrm{Ryd}}=3.28984\times10^{15}$\,Hz).
$J_\nu$ is the angle-averaged specific intensity (the mean intensity), so that
$u_\nu=4\pi J_\nu/c$ is the spectral energy density.

\begin{longtable}{@{}L{0.22\linewidth}L{0.73\linewidth}@{}}
\caption{Notation and units.}\label{tab:notation}\\
\toprule
Symbol & Meaning \\
\midrule
\endfirsthead
\toprule
Symbol & Meaning \\
\midrule
\endhead
$n_{\mathrm H}$ & density of hydrogen nuclei, \rtn{grid\%Hden} [cm$^{-3}$] \\
$x_i = n_i/n_{\mathrm H}$ & fractional abundance of species $i$, \rtn{grid\%pdrAbun} [dimensionless] \\
$J_\nu$ & mean intensity [erg\,s$^{-1}$\,cm$^{-2}$\,Hz$^{-1}$\,sr$^{-1}$] \\
$k$, $\nu_k$, $\Delta\nu_k$ & index, tabulated frequency and width of a bin of the MC frequency grid
            (\rtn{nuArray}, \rtn{widFlx}) \\
$\Jste_k$ & MC estimator of the \emph{total} field in bin $k$,
            $4\pi\int_{\mathrm{bin}\,k} J_\nu\,\dd\nu$ [erg\,s$^{-1}$\,cm$^{-2}$], \rtn{grid\%Jste} \\
$\Jdif_k$ & the part of $\Jste_k$ carried by scattered or re-emitted packets, \rtn{grid\%Jdif} [same units] \\
$\uD$ & $8.94\times10^{-14}$\,erg\,cm$^{-3}$, the 6--13.6\,eV energy density of the
        \citet{Draine1978} interstellar field (\src{constants\_mod.f90:\allowbreak{}29});
        $c\,\uD=2.68\times10^{-3}$\,erg\,s$^{-1}$\,cm$^{-2}$ \\
$\chi$ & band-integrated (6--13.6\,eV) field in Draine units, a dimensionless energy density,
         $\chi = \sum_{k\in\mathrm{FUV}} \Jste_k/(c\,\uD)$; $\chi=1$ is $G_0=1.69$ in
         Habing units of the same band (Section~\ref{sec:bands}) \\
$\chi_0$ & geometrically diluted, unattenuated field of the sources, same units (Equation~\ref{eq:chi0}) \\
$\chiD,\ \chiF$ & direct (never scattered or re-emitted) and diffuse parts of $\chi$ \\
$s_{\mathrm D}$ & share of the Draine field's 6--13.6\,eV energy in the Lyman--Werner (LW) window 11.2--13.6\,eV \\
$f_{\mathrm{dir}},\ f_{\mathrm{dif}}$ & LW window factors of the direct and diffuse parts (Section~\ref{sec:window}) \\
$\eta_{lu}$ & photon occupation number at the line $l\to u$, $J_\nu c^2/(2h\nu^3)$ [dimensionless] \\
$\kappa^{\mathrm{abs}}_\nu,\ \kappa^{\mathrm{sca}}_\nu$ & dust absorption and scattering coefficients
         [cm$^{-1}$], \rtn{absOpac}, \rtn{scaOpac}; $\kappa_V$ their sum at 5500\,\AA \\
$\AV$ & visual extinction [mag]; $1.0857 = 2.5\log_{10}e$ mag per unit optical depth (\rtn{magPerTau}) \\
$\RFUV$ & ratio of the band-mean FUV dust extinction coefficient to $\kappa_V$, per cell (Section~\ref{sec:dust}) \\
$N_X$ & column density of species $X$ [cm$^{-2}$]; $N_l$ the column of \Htwo{} in level $l$ \\
$b$ & Doppler parameter, $b^2 = 2kT/m_{\mathrm{H_2}} + v_{\mathrm{turb}}^2$ [cm\,s$^{-1}$] \\
$\tau_0,\ a$ & line-centre optical depth and Voigt damping parameter of an \Htwo{} line \\
$S(\tau_0,a)$, $W(\tau_0,a)$ & single-line self-shielding function (profile-weighted transmission of a continuum in a Voigt line) and its integral over $\tau_0$ (curve of growth) \\
$k_{\mathrm{Sol}}$ & Solomon photodissociation rate per \Htwo{} molecule [s$^{-1}$], \rtn{grid\%h2Sol} \\
$R_{\mathrm{f}}$ & \Htwo{} grain formation coefficient [cm$^3$\,s$^{-1}$]; formation rate per volume $R_{\mathrm f}\,n_{\mathrm H}\,n(\mathrm H)$ [cm$^{-3}$\,s$^{-1}$] \\
$\zeta$ & cosmic-ray ionisation rate, \kw{pdrZeta} [s$^{-1}$] \\
$T_e$, $T_d$ & gas (electron) temperature and mean grain temperature [K]; $T_{\mathrm{CMB}}=2.725$\,K \\
\bottomrule
\end{longtable}

\subsection{Who owns what}
\label{sec:ownership}

The coupling is designed so that every quantity has a single writer
(\src{CLAUDE.md}, ``one reality per quantity''). Table~\ref{tab:owner}
lists the quantities that lie on the path from the field to the molecules.

\begin{longtable}{@{}L{0.26\linewidth}L{0.16\linewidth}L{0.25\linewidth}L{0.24\linewidth}@{}}
\caption{Ownership of the quantities on the path from the radiation field to the molecular content.}\label{tab:owner}\\
\toprule
Quantity & Owner & Written in & Read by \\
\midrule
\endfirsthead
\toprule
Quantity & Owner & Written in & Read by \\
\midrule
\endhead
$J_\nu$ in every bin (total and diffuse) & MOCASSIN MC & \rtn{pathSegment}, \src{photon\_mod.f90:\allowbreak{}1900--1924, 2236--2250} & everything below \\
dust opacity, albedo, $T_d$ & MOCASSIN & \src{dust\_mod.f90}, \rtn{getDustT} (\src{update\_mod.f90}) & $\RFUV$, \Htwo{} formation, level model \\
$T_e$ & MOCASSIN \rtn{thermBalance} & \src{update\_mod.f90:\allowbreak{}1734} & chemistry, level model \\
ionisation stages of H, He, C, O, Mg (and all other elements), $n_e$ & MOCASSIN \rtn{ionBalance} & \src{update\_mod.f90:\allowbreak{}2527} & imposed on the chemistry \\
atom photoionisation and recombination rates of H, He, C, O, Mg; \Htwo{}$\to$\Htwo$^+$ & MOCASSIN & \src{update\_mod.f90:\allowbreak{}2840--2878} & chemistry (replace network channels) \\
$\AV$ and the source-ward columns $N(\mathrm{H_2}), N(\mathrm C), N(\mathrm{CO})$ & bridge geometry & \rtn{computeAVradial}, \src{pdr\_mod.f90:\allowbreak{}1859} & chemistry \\
scattering sites and their two-leg columns & bridge geometry & \rtn{traceDiffuseSites}, \src{pdr\_mod.f90:\allowbreak{}2136} & chemistry, level model \\
\Htwo{} level populations, $k_{\mathrm{Sol}}$, \Htwo{} heating & level model & \rtn{h2LevelModel}, \src{pdr\_mod.f90:\allowbreak{}5201} & chemistry, thermal balance \\
abundances of the 33 species & chemistry & \rtn{pdrStep}, \src{pdr\_mod.f90:\allowbreak{}3180} & MOCASSIN, level model, columns \\
fraction of each element not in molecules (\rtn{elemAvail}) & chemistry & \rtn{setElemAvail}, \src{pdr\_mod.f90:\allowbreak{}5120} & \rtn{ionBalance} \\
$n(\mathrm{H_2})$ for the gas opacity (\rtn{H2den}) & chemistry & \src{pdr\_mod.f90:\allowbreak{}4573} & \rtn{addOpacity} \\
cooling by CO & chemistry (3D-PDR level populations, tabulated in $T$) & \rtn{pdrStep} & \rtn{thermBalance} \\
cooling by \Cp, C, O & MOCASSIN's own atoms, with the chemistry's neutral colliders and escape probabilities (\dlog{D-363}) & \src{emission\_mod.f90:\allowbreak{}1861--1887} & \rtn{thermBalance}, line fluxes \\
heating terms MOCASSIN lacks (\rtn{pdrHeatXtra}) & chemistry and level model & \src{pdr\_mod.f90:\allowbreak{}4499--4546} & \rtn{thermBalance} \\
\bottomrule
\end{longtable}

\subsection{Overview of the chain}

Figure~\ref{fig:chain} shows the chain inside one MC iteration. The field
reaches the molecules through five channels:
(i) directly, through the photodissociation and photoionisation rates, which
read the local MC field (Sections~\ref{sec:field} and \ref{sec:photorates});
(ii) through the gas temperature, which sets the rate coefficients, the
collisional excitation of \Htwo{} and, through the thermal speed and the
sticking of H atoms, the \Htwo{} formation rate;
(iii) through the dust temperature, which sets the efficiency of \Htwo{}
formation on grains (Section~\ref{sec:Td});
(iv) through MOCASSIN's ionisation balance, which fixes the \Cp/C ratio and
the electron density the chemistry works with; and
(v) through the molecules themselves, whose columns, integrated along the
paths the photons took, shield the \Htwo{} and CO lines at the next
iteration (Section~\ref{sec:geometry}).
The molecules act back on the transfer only through the \Htwo{} photoionisation
continuum and, indirectly, through the ionisation and energy balances
(Section~\ref{sec:feedback}); no molecular line opacity enters the MC transfer.

\begin{figure}[tbp]
\centering
\small
\fbox{\parbox{0.86\linewidth}{\centering
\textbf{1. MC transfer} (\rtn{energyPacketDriver}): $\Jste_k$, $\Jdif_k$ in every cell and bin;
dust temperatures; packets flagged as never scattered (direct) or scattered/re-emitted (diffuse)}}\\[2pt]
$\Big\downarrow$\\[2pt]
\fbox{\parbox{0.86\linewidth}{\centering
\textbf{2. Dust and geometry} (\rtn{dustPerCell}, \rtn{computeAVdust}):
$\RFUV$, $\AV/N_{\mathrm H}$, grain area per H;
exact rays to the sources: $\AV$ and $N(\mathrm{H_2})$, $N(\mathrm C)$, $N(\mathrm{CO})$ of iteration $n-1$'s abundances;
snapshot of the scattering emissivity for the diffuse walk; $\chi$, $\chi_0$}}\\[2pt]
$\Big\downarrow$\\[2pt]
\fbox{\parbox{0.86\linewidth}{\centering
\textbf{3. \Htwo{} level model} (\rtn{h2LevelModel}, one sweep): per line
$\eta = \eta_{\mathrm{dir}}\langle S\rangle_{\mathrm{dir}}\mathcal O_{\mathrm{dir}}+\eta_{\mathrm{dif}}\langle S\rangle_{\mathrm{dif}}$;
Voigt self-shielding with level columns, line overlap, trapping;
302-level statistical equilibrium at $T_e^{(n-1)}$ $\Rightarrow$ $k_{\mathrm{Sol}}$, \Htwo{} heating}}\\[2pt]
$\Big\downarrow$\\[2pt]
\fbox{\parbox{0.86\linewidth}{\centering
\textbf{4. MOCASSIN cell update} (\rtn{updateCell}): ionisation and $T_e$ of every cell,
with the chemistry's channels of iteration $n-1$;
stores the photoionisation and recombination rates of H, He, C, O, Mg and the \Htwo{} photoionisation rate}}\\[2pt]
$\Big\downarrow$\\[2pt]
\fbox{\parbox{0.86\linewidth}{\centering
\textbf{5. Chemistry} (\rtn{pdrStep}): rays $=$ direct beam $+$ scattering sites, all at $\AV=0$;
generic photorates $\alpha\chi$; CO with \citet{vanDishoeckBlack1988} shielding and window factor;
\Htwo{} photodissociation $=k_{\mathrm{Sol}}$; MOCASSIN's ionisation imposed;
CVODE to steady state; CO cooling table; heating terms}}\\[2pt]
$\Big\downarrow$\\[2pt]
\fbox{\parbox{0.86\linewidth}{\centering
\textbf{6. Feedback} (\rtn{pdrToMocassin}, \rtn{rescaleIonToAvail}): fraction of each element in molecules
$\rightarrow$ ion stages; $n(\mathrm{H_2})\rightarrow$ gas opacity; coolants and heating $\rightarrow$ next $T_e$;
new abundances $\rightarrow$ next iteration's shielding columns; second ionisation solve (kept at the end)}}
\caption{The chain from the radiation field to the molecular content inside one MC iteration $n$ of the
\texttt{h2-model} branch (3D build), in the order of execution. Inputs marked ``iteration $n-1$'' or
with the superscript $(n-1)$ are one iteration old (Section~\ref{sec:lags}); the lags vanish at a fixed
point of the iteration. The steps and their line numbers are listed in Table~\ref{tab:steps}.}
\label{fig:chain}
\end{figure}
